\chapter{Diffraction, Interference, and Polarization}
\label{chap:diffraction-interference-polarization}

The previous chapter read a field's holonomy and watched the instrument prove its first loop residue. This
chapter reads what a field does when it overlaps with itself --- when it diffracts past an obstacle,
interferes with its own other path, and is filtered by orientation --- and in doing so it introduces two of
the instrument's recurring moves. One is the \emph{distinguishability test}: the instrument shows
interference exactly when it cannot tell which path a record took, and the fringes are its own ignorance made
visible. The other is the chapter's wall --- the second no-go, where crossed filters pass nothing, an
impossibility and not a small number. Four effects carry it: a bright spot exactly where ray optics demands
darkness, two ledgers a fold reconciles, the wall of crossed polarizers, and the double slit, where the
record's ability to tell the paths apart, and nothing else, decides whether a pattern survives.

\section{The Arago (Poisson) Spot: symmetry forbids darkness}
\label{se:arago}

Poisson offered the prediction as a reductio against the wave theory: if light were a wave, there would be a
bright spot at the very center of the shadow cast by a circular disk --- a place ray optics insisted must be
perfectly dark. Arago performed the experiment and found the spot. To the instrument the spot is not a
surprise but a necessity, and the necessity is a symmetry. Every point on the disk's circular rim lies at
exactly the same distance from the axial point directly behind the disk, so every secondary wavelet the rim
emits arrives at that point having travelled the same way --- all of them in phase, none of them out. The
instrument, tallying what arrives on axis, is summing a ring of contributions that are all aligned,
\begin{equation}
  A_{\mathrm{axial}} = \oint_{\mathrm{rim}} dA \quad (\text{all in phase}) \;\neq\; 0,
\end{equation}
and this integral is not a Fresnel computation to be ground out; it is the \emph{name} of an in-phase sum
that cannot vanish. For the center to be dark, the contributions around the rim would have to cancel, and
cancellation needs a phase asymmetry --- some points arriving ahead, others behind. The perfect circular
symmetry of the rim forbids exactly that asymmetry. There is no distinction around the rim for the
cancellation to use, and so the instrument cannot make the center dark.

The honest floor is a finite count obligation: the bright spot is forced by the symmetric, in-phase count
around the rim, and a dark center would demand an asymmetry that is simply not in the record. The reading is
that the spot is not waves somehow bending into the shadow but symmetry forbidding darkness. Where the
symmetry is exact, the cancellation a dark center needs has nothing to grip, and what cannot be unrecorded is
bright. Poisson's reductio was right about the spot and wrong about the absurdity: the wave theory does demand
it, because the instrument's own symmetry demands it.

\section{The Mach--Zehnder Effect: two ledgers until a fold}
\label{se:mach-zehnder}

The interferometer shows the instrument keeping two ledgers at once and reconciling them at a boundary. A
single photon enters a Mach--Zehnder interferometer, and at the first beam splitter the record divides into
two coexisting, consistent ledgers --- the two arms. Each accumulates its own refinements --- reflections,
delays, phase shifts --- recorded as an accumulated phase, and no experiment performed within a single arm
can determine which arm it is. The instrument genuinely carries both; neither is the real one and the other a
ghost. At the second beam splitter --- the fold --- the two ledgers are brought to a shared edge and
reconciled,
\begin{equation}
  U_{\mathrm{final}} = f\bigl(U^{(1)}, U^{(2)}\bigr),
\end{equation}
and the accumulated phase difference, taken modulo $2\pi$, decides what the reconciled record shows. Written as
amplitudes, the reconciliation is a sum. Each arm carries an amplitude, $\psi_1$ and $\psi_2$, and the fold
adds them,
\begin{equation}
  \psi = \psi_1 + \psi_2,
\end{equation}
so the intensity the instrument reads is the squared magnitude of the sum,
\begin{equation}
  I = |\psi_1 + \psi_2|^2 = I_1 + I_2 + 2\sqrt{I_1 I_2}\,\cos\Delta\varphi,
\end{equation}
the two arm intensities plus a cross-term. That cross-term, $2\sqrt{I_1 I_2}\cos\Delta\varphi$, is the whole
of the interference: positive when the phases agree ($\Delta\varphi = 0$, bright), negative when they oppose
($\Delta\varphi = \pi$, dark), and exactly what a plain sum of two independent ledgers could never produce.
Interference is the cross-term, and the cross-term survives only because the two ledgers were reconciled as
amplitudes at the fold, not added as separate intensities. If the two phases agree modulo $2\pi$ the ledgers
merge without creating new distinctions and the output is bright; if they differ by half a turn they cancel
and the output is dark. The interference law $I \propto
\cos^2(\Delta\varphi/2)$ is the \emph{shape} of this reconciliation, the smooth contour the fringe traces as
the phase difference is varied --- but the act underneath it is the instrument's boundary reconciliation, the
same fold-at-an-edge it ran to fix the GPS interior in Part I.

The honest floor is a finite count obligation: which arm carries the count, and how the two ledgers reconcile
at the fold, is a boundary question the instrument settles by the phase the record carried around. The
reading is that the two paths are not one taken and one forgone but two coexisting ledgers a fold reconciles.
Interference is that reconciliation --- bright or dark according to the phase difference the two ledgers bring
to their shared edge --- and it is the same boundary reconciliation that will, in the quantum part, decide a
measurement's outcome from the frontier where two records meet.

\section{The Malus Effect: the wall of crossed filters}
\label{se:malus}

Here the instrument meets its second wall, and it is the same kind of wall it met at static friction. A beam
that has passed a linear polarizer meets a second polarizer set at an angle $\theta$ to the first, and the
fraction that survives is
\begin{equation}
  I = I_0 \cos^2\theta .
\end{equation}
Read $\cos^2\theta$ as a \emph{shape}, not a number to compute: it is the profile of how much of the beam's
orientation projects onto the second filter's axis --- full at $\theta = 0$, where the axes align, tapering
as the angle opens. The projection has a precise form. Write the beam's polarization as a unit vector --- the
Jones vector $\hat a$ --- pointing along its axis of oscillation, and the analyzer's axis as a second unit
vector $\hat e$; the fraction of intensity that survives is the squared projection of one onto the other,
\begin{equation}
  I = I_0\,|\hat a \cdot \hat e|^2 = I_0 \cos^2\theta,
\end{equation}
since $\hat a \cdot \hat e = \cos\theta$ for two unit vectors at angle $\theta$. Malus's law is geometry: the
surviving amplitude is the component of the polarization along the analyzer, and the surviving intensity its
square. And at $\theta = 90^\circ$, the filters crossed, the projection is zero. Not small:
zero. No photon passes; the record ends; the light can no longer be observed. The honest floor here is not a
count that balances but an impossibility. At the crossed angle, transmission is inadmissible --- there is no
refinement of the record that produces a transmitted photon the law would permit, exactly as no refinement
produced a sub-threshold slip at the friction wall. As there, the law is honest about its character: it does
not claim each photon is absorbed and re-emitted, only that, crossed, the beam does not pass.

Watch the instrument turn the second filter. At $\theta = 0$ the count comes through in full; turn the filter
a little and the count drops, following the $\cos^2$ contour down; keep turning and it dwindles toward the
crossed angle. At $90^\circ$ it reaches not a small residue but a flat zero, and there it stays --- no jitter,
no leak, no refinement that recovers a single transmitted photon. The instrument can sit at the wall and read
the record as finely as it likes; the wall does not soften. That is what marks a no-go apart from a small
count: a small count shrinks as you look closer, and a wall does not move.

The honest floor is a no-go, the second of the four. The reading is that crossed polarizers extinguish a beam
by forbidding its transmission, not by accounting for each photon's fate. The $\cos^2\theta$ shape describes
the surviving projection at every other angle; the wall is the one direction, $\theta = 90^\circ$, where the
projection vanishes and the instrument can do nothing but record that nothing comes through. The wall recurs
throughout the book --- a domain seam that cannot be misdrawn, a closed loop in time the count's own rising
forbids --- and it recurs as exactly this: a refinement the law makes impossible.

\section{The Schr\"odinger--Young Effect: which-path destroys the pattern}
\label{se:schrodinger-young}

The double slit makes the deepest point in the simplest geometry, and it is the clearest statement of the
instrument's distinguishability test. With both slits open and nothing in the record able to tell which slit a
particle passed, the two amplitudes add coherently, and the intensity carries an interference term,
\begin{equation}
  |\psi_A + \psi_B|^2 = |\psi_A|^2 + |\psi_B|^2 + 2\,\mathrm{Re}\bigl(\psi_A^{*}\psi_B\bigr),
\end{equation}
the cross term being the fringe pattern itself. Now introduce which-path information --- anything that lets
the record distinguish the two slits, a tag on the particle, a recoil in the screen --- and the cross term
vanishes; the intensity collapses to the plain sum,
\begin{equation}
  |\psi_A|^2 + |\psi_B|^2,
\end{equation}
with no fringes at all, even when no continuous trajectory was ever recorded. The decisive quantity is not
whether anyone looked but whether the record \emph{can} tell the paths apart.

This is the instrument's distinguishability test, and it is worth watching directly. Set up the slits and
turn the which-path detector off: the fringes appear, sharp and full. Turn the detector on, so the record now
carries which slit each particle took: the fringes vanish, replaced by the featureless sum. Turn it off
again: the fringes return. Nothing about the particles' paths has changed; what changed is whether the record
could distinguish them. The visibility of the fringes is, exactly, a measure of the instrument's own
which-path ignorance --- full fringes when it cannot tell, none when it can, and a smooth trade between. The
pattern is not destroyed by observation; it is destroyed by distinguishability.

The trade has an exact form, and it is one of the sharpest statements in the subject. Measure the contrast of
the fringes by their visibility,
\begin{equation}
  V = \frac{I_{\max} - I_{\min}}{I_{\max} + I_{\min}},
\end{equation}
which is one for perfect fringes and zero for a flat sum. Measure how well the record can tell the paths apart
by a distinguishability $D$, zero when the paths are indistinguishable and one when the record fixes the slit.
The two are bound by a single inequality,
\begin{equation}
  V^2 + D^2 \le 1,
\end{equation}
the wave--particle duality relation. It is not a verbal slogan but a quantitative ceiling: full fringes
($V = 1$) demand total ignorance of the path ($D = 0$); full path knowledge ($D = 1$) forbids any fringe
($V = 0$); and in between the instrument can trade one for the other but never exceed the bound. The pattern's
visibility and the record's which-path knowledge are two readings of one quantity, and their squares cannot
together pass one.

The honest floor is a finite count obligation: the interference term survives exactly when the paths are not
distinguished, so the obligation is on the distinguishability, not on a trajectory. The reading is that the
pattern is the instrument's ignorance made visible. What erases the fringes is the record's capacity to
separate the paths, and that capacity --- not any act of looking --- is the physical content. The effect
returns as a central character of the quantum part, where the same distinguishability decides not only fringes
but the outcome of a measurement; read here, among the fields, it is wave interference, and the lesson it
teaches is the one the quantum will build on: distinguishability is the variable that matters.

\bigskip

Diffraction and interference have shown the instrument two things it will use everywhere. A pattern of light
is a reconciliation of ledgers --- forced bright by a symmetry the record cannot break, or bright-and-dark by
a phase two arms bring to a fold --- and it survives exactly insofar as the record cannot tell the paths
apart. And crossed filters have given the instrument its second wall, an angle where transmission is simply
not admissible. The next chapter stays with light but turns to the limits of reading it: how finely a wave can
be sampled, what a finite spectrum does to a sharp edge, and what a measurement can and cannot recover from
the counts it takes.
